\setcounter{section}{56} % tema número N-1
\section{Multi Usuario}

Nombre completo: Tipos de sistemas de información multiusuario. Sistemas grandes, medios y pequeños. Servidores de datos y aplicaciones. Virtualización de servidores.

\localtableofcontents

\subsection{Sistemas de información multiusuario}
    La clasificación es aproximada, y tiene herencias legacy. De mayor a menor son

    \begin{description}
        \item[Mainframes] Soportan una elevada carga de trabajo y usuarios. Son estructuras centralizadas. Tienen una alta fiabilidad y se diseñan para trabajar sin interrupción durante años
        \item[Super ordenadores] Gran potencia en base a arquitecturas paralelas. Se gestionan por personal especializado.
        \item[Miniordenador o Servidor] Ordenador medio. Puede realizar cálculos complejos y soporta una gran cantidad de usuarios. Están a mitad de camino entre el mainframe y el microordenador.
        \item[Estaciones de trabajo] Ordenadores personales con altas prestaciones. Se usan comúnmente para CAD o diseño 3D
        \item[Microordenadores] Ordenadores de sobremesa o portátiles.
        \item[Otros] Tablets y smartphones.
    \end{description}

\subsection{Servidores de datos y aplicaciones}
    Según el papel que asumen se pueden clasificar como \textbf{dedicados} si se usan para un único fin. o \textbf{no dedicados} si aceptan varios servicios simultáneamente.

    Los usos mas comunes son
    \begin{description}
        \item[web] se encargan de la capa de presentación
        \item[aplicaciones] Ejecutan la lógica de negocio
        \item[datos] Sistemas de base de datos
        \item[correo]
        \item[multimedia] Distribuye contenido multimedia
        \item[noticias]
        \item[colaboración] Provee de herramientas colavorativas para sincronizar el trabajo de varios usuarios
        \item[proxy] Intermediario de conexión para sus usuarios
    \end{description}

\subsection{Virtualización}
    Gracias a tecnologías como la nube se permite una reducción de costes operativos. Cambiando el paradigma a los servicios por uso en lugar de por disponibiilidad. Gracias a ello se pueden hacer despliegues automáticos, escalado bajo demanda o continuidad de negocio.

    Según cuanto se virtualiza y qué capa queda expuesta, la virtualización puede ser de varios tipos.
    \begin{description}
        \item[Emulación] se emulta completamente un tipo de arquitectura en otra. (DosBox)
        \item[Nativa o total] La máquina virtual emula una cantidad suficiente de hardware para que el entorno virtualizado funcione completamente sin modificar (VMWare, VirtualBox)
        \item[para-virtualización] El host no emula el hardware, sino que ofrece una serie de APIs a un sistema convenientemente modificado (UML, xVM)
        \item[Virtualización a nivel de SO] Los huéspedes comparten el mismo kernel que el host. Creandose instancias independientes entre si. (Linux V Server, OpenVZ)
        \item[Virtualización de aplicaciones] Las aplicaciones tienen su propio entorno virtualizado con lo necesario para ejecutarse tanto sobre un servidor, como un cliente.
    \end{description}

\subsection{Generaciones de Ordenadores}
    Para hablar de tamaños de ordenadores se usan escalas dependientes del número de transistores, repartidos en órdenes de magnitud:
    \begin{table}[H]
    \begin{tabular}{ll}
        SSI & Small Scale Integration \\
        MSI & Medium Scale Integratrion \\
        LSI & Large Scale Integration\\
        VLSI & Very Large Scale Integration\\
        ULSI & Ultra Large Scale Integration\\
    \end{tabular}
    \end{table}

    \subsubsection{1940 a 1956}
        El Eniac es el primer ordenador de propósito general. Establecido en la universidad de Pensilvania, ocupaba \begin{math}167m^2\end{math}. Utilizaba válvuilas de vacío, utilizaba 160KW y generaba 50° en la sala. Von Neuman establece las referencias de lo que debe ser una máquina multipropósito.

        El EDVAC es el primer ordenador con memoria. Utiliza diodos y tubos de vacío. La programación se hacía físicamente y usaba un lenguaje máquina.

    \subsubsection{1955 a 1964 (SSI)}
        Uso de transistores y circuitos integrados. Programación a mas alto nivel con Cobol y Fortran. Ordenadores como el UNIVAC y el IBM.

    \subsubsection{1964 a 1971 (MSI - LSI)}
        Primeros microprocesadores gestionando sistemas enteros como el 4004 y el 8008. Empiezan los lenguajes de alto nivel como Pascal o Basic. Las unidades tienen un alto coste, asumible por grandes corporaciones. Empiezan los primeros Mainframes como IBM S/360, Cray o ILLIAC

    \subsubsection{1977 a 1989 VLSI}
        Principio de la informática personal con tamaño reducido. Ordenadores como el Apple I y II, IBM PC o el Macintosh.

    \subsubsection{1980 a Actualidad}
        No ha cambiado la forma, solo la potencia y el tamaño. Ordenadores multiprocesador, multiprocesamiento, distribución, conexiones de red y popularización.
\subsection{Componentes de sistemas multiusuario}
    Un sistema es multiusuario si da servicio compartiendo recursos simultáneanos a múltiples usuarios concurrentes. Según sus prestaciones son pequeños, medianos o grandes. No se ponen cifras por su tamaño.
    
    Los sistemas grandes son o Mainframes o Supercomputadores. Un mainframe se usa para carga pesada y crítica. Soporta carga con mucho IO, almacenamiento, etc. Todo ocurre en el mismo dispositivo. Tiene un alto coste de licencia, fabricante, certificación\dots Puede generar clientes cautivos. También es poco popular en favor de los supercomputadores. Un supercomputador por otro lado está pensado para cargas paralelas, menos complegidad y ES.

    Los sistemas pequeños atienden a un número limitado de usuarios. Se usan en entornos reducidos. Se asemejan a ordenadores personales o blades.

    Los sistemas medios están en la frontera entre ambos. Para dimensionar un sistema se suele tener en cuenta el número de usuarios, la carga, disponibilidad, aplicaciones y los criterios del ENS.

    Otras clasificaciones tienen en cuenta la utilidad, interoperabilidad con otros sistemas o la arquitectura. También puede haber requisitos en cuanto a la arquitectura.

\subsection{Servidores}
Según la RFC 2616 sobre HTTP, es un software que atiende peticiones. Se denomina bundle al conjunto de elementos que permiten atender estas peticiones. 

A la hora de dividir, se suele hacer entre cliente y servidor. Aunque hay mas divisiones subyacentes, como la arquitectura de 3 capas, que divide el cliente en cliente y presentación y el servidor entre el Software de Aplicación y el Software de datos. Ulteriormente, en N capas, se divide la aplicación en diferentes applets. La de presentación se encarga de servir de interfaz entre la parte pública y privada, filtrando y delegando a las capas siguientes, que quedan aisladas del exterior.